\section{Exploring the Cloud Landscape}
\label{section:cloud}

\newcommand{\pie}[1]{%
\begin{tikzpicture}
 \draw (0,0) circle (1ex);\fill (1ex,0) arc (0:#1:1ex) -- (0,0) -- cycle;
\end{tikzpicture}%
}

\newcommand{\revpie}{%
\begin{tikzpicture}
 \draw (0,0) circle (1ex);\fill (0,1ex) arc (90:270:1ex) -- (0,0) -- cycle;
\end{tikzpicture}%
}

\newcommand{\tableYes}{\pie{360}}
\newcommand{\tableNo}{\pie{0}}
\newcommand{\tablePartial}{\revpie}

\begin{table}[!t]
    \renewcommand{\arraystretch}{1.35}
    \newcolumntype{L}{>{\raggedleft\arraybackslash}X}
    \newcolumntype{C}{>{\centering\arraybackslash}X}
    \newcolumntype{R}{>{\raggedright\arraybackslash}X}
    \centering

    \begin{tabularx}{.9\linewidth}{l C C C}
        \toprule
                                           & AWS           & Azure         & GCP           \\
        \hline                                                              
        TEE                                & \tableYes     & \tablePartial & \tableYes     \\
        Firmware                           & \tableYes     & \tablePartial & \tablePartial \\
        Kernel                             & \tablePartial & \tablePartial & \tablePartial \\
        Root FS                            & \tablePartial & \tablePartial & \tablePartial \\
        \hline                                                              
        \textbf{Nominal AL}                & \textbf{4}    & \textbf{4}    & \textbf{4}    \\
        \textbf{Trustworthy AL}            & \textbf{2}    & \textbf{0}    & \textbf{1}    \\
        \bottomrule
    \end{tabularx}

        \vspace{0.2cm}
        {\raggedright \tableYes\xspace= Verifiable w/o trust; \tablePartial\xspace= Verifiable w/ trust; \tableNo\xspace= Not verifiable \par}
        \vspace{0.2cm}
        \caption{AMD \sevsnp{} offerings on \ac{AWS}, Azure and \ac{GCP}. This
          table reflects attestation levels described in
          \cref{section:att-levels} and verifiability with or
          without trust, resulting in the maximum nominal AL (with trust in the
          \ac{CSP}) and trustworthy AL (without trust). Results dated \clouddate.}
        \label{tbl:cloud-landscape}
        %\vspace{-5mm}
\end{table}
%
  


We investigated current \ac{CC} offerings on public clouds, focusing on
\sevsnp{} as the most widely available VM-based \ac{TEE}. Our evaluation is
solely based on commercial cloud features under general availability or public
preview. We found that only the three major \acp{CSP}, i.e., Microsoft Azure,
\ac{GCP} and \ac{AWS}, offer to deploy \sevsnp{} \acp{CVM} on demand. We looked
for VM-based confidential computing support in other \acp{CSP}, but they either
supported plain \sev{} or only offered single-tenant bare-metal servers with
\sevsnp{} hardware. While plain \sev{} does not provide integrity protection nor
flexible attestation capabilities, bare-metal servers come at much higher costs
and typically require a fixed monthly subscription. Instead, we were interested
in the multi-tenant scenario where customers can deploy \acp{CVM} on demand on a
shared infrastructure, with little to no control over the hypervisor.

As discussed in \cref{sec:scenario}, our evaluation tried to determine the level
of \emph{verifiability without trust} that can be achieved on such
\acp{CSP}. Besides, we also considered the model where the verifier (partially)
trusts the \ac{CSP} and relies on provided services and data, such as a custom
attestation report or a proprietary attestation service.

Results are shown in \cref{tbl:cloud-landscape} and based both on public
documentation and direct experiments. Regarding the former, all references in
\cref{eval:cloud-results}, except Github links, have been uploaded to the
Internet Archive\footnote{\url{https://www.archive.org}} to provide a stable
reference. Concerning the latter, our evaluation was made on \clouddate{}. Since
this is a snapshot of the current state and features of clouds, some of our
findings may not hold in the future. However, our evaluation framework is still
useful to evaluate future offers in the \ac{CVM} area.

\subsection{Results}
\label{eval:cloud-results}

In our experiments on \ac{AWS} and \ac{GCP}, we were able to fetch a raw
attestation report from the AMD \ac{SP} via the \sevguest{} device, allowing to
independently verify that the guest VM is indeed running in a up-to-date
\sevsnp{} \ac{TEE}. Azure \acp{CVM}, instead, only provide access to an
attestation report generated at boot time and stored in \ac{vTPM} non-volatile
memory. As this report does not contain any freshness information chosen by the
verifier, the latter cannot distinguish a legitimate attestation report from a
replay attack and must rely on the \ac{CSP} to attest that the VM is indeed
using \sevsnp. We tried to boot an \snp{}-aware kernel built from the AMDSEV
repository~\cite{amdsevGithub}, yet we still could not access the AMD \ac{SP}.

Regarding AL2, only \ac{AWS} offers an open source firmware that can be
reproducibly built~\cite{awsUefi}. We successfully checked that the attestation
report indeed reflects the correct measurement. Instead, Azure and \ac{GCP}
guests boot a proprietary firmware whose code cannot be audited. To the best of
our knowledge, firmware reference measurements are also not provided, meaning
that the verifier cannot even compare the measurement in the attestation report
against the expected value. However, we observed that both Azure and \ac{GCP}
have a secure boot option that prevents \acp{CVM} from booting if the firmware
is corrupted~\cite{azureDocsSecureBoot,googleDocsSecureBoot}, though we could
obviously not experimentally verify this claim. Yet, tenants that are willing to
trust their \ac{CSP} can be confident that a running firmware was not tampered
with by an outsider attacker.

All \acp{CSP} allow extending the measurements up to the kernel via a \ac{vTPM}.
However, on \ac{AWS}~\cite{awsGithubvTPM} and \ac{GCP}~\cite{googleDocsvTPM} the
\ac{vTPM} is implemented in the hypervisor, greatly increasing the \ac{TCB}.
Azure, instead, exposes a \ac{vTPM} from within the \ac{CVM}
firmware~\cite{azureDocsvTPM} which, as mentioned before, cannot be reviewed.
Thus, no \ac{CSP} currently supports trustworthy attestation of kernel
components.

After kernel verification, the chain of measurements can be extended to the root
filesystem via different approaches (\cref{section:att-examples}).  Hence, AL4
is a natural extension of AL3, provided that the \ac{CSP} allows for the
customization or upload of VM images by the guest owner. We experimentally
verified that the former is indeed possible by installing a custom bootloader
and kernel, as well as changing kernel command-line parameters, e.g., to enable
Linux \ac{IMA}. Uploading images seems also feasible according to public
documentation~\cite{awsDocsCustomImage,azureDocsCustomImage,googleDocsCustomImage},
but we did not verify this. Still, for trustworthy attestation, the integrity of
the root filesystem ultimately relies on the integrity of the boot chain up to
the kernel: If the latter cannot be verified independently of the \ac{CSP},
neither can the former.

In summary, all \acp{CSP} provide the infrastructure to perform a full
attestation (\emph{nominal AL4}), but they rely on either proprietary software
or components outside the \ac{TEE} TCB, or both. On Azure, a verifier can obtain
an attestation report but cannot verify its freshness, hence replay attacks are
possible and the \ac{CSP} must be trusted for all attestation levels
(\emph{trustworthy AL0}). On \ac{GCP}, one can fetch dynamic reports, but the
firmware cannot be independently verified (\emph{trustworthy AL1}). Only
\ac{AWS} allows to attest an auditable firmware, yet extending the chain of
measurements to later boot stages uses a hypervisor-based \ac{vTPM} outside the
\ac{TEE} (\emph{trustworthy AL2}).

\subsection{Discussion}

We showed that guest owners still require a significant amount of trust
in the \acp{CSP} when using \ac{CVM} solutions in the cloud. This goes against
the classic \ac{TEE} threat model, where the infrastructure is considered
untrusted. There are two reasons for this:
%
\begin{inparaenum}
    \item \ac{CSP} code is running inside the \ac{TEE} boundary, with no
    possibility for review by the remote verifier, and
    \item the attestation process involves \ac{CSP}-managed components.
\end{inparaenum}
%
While running \ac{CSP} code in the \ac{CVM} is sometimes necessary (e.g., to
configure network interfaces, mount devices, etc.), it is certainly possible to
increase transparency towards the verifier, as demonstrated by \ac{AWS}'
reviewable firmware. The second problem is mostly caused by a lack of
customization options offered by \acp{CSP}. This is somewhat understandable
since most cloud customers are looking for ease of use and ``one-click
compliance'', where they are willing to trust the \ac{CSP} to perform security
checks in their place. For advanced users, however, more customizability is
desirable.

%There is also a fundamental challenge with VM-based \acp{TEE}, which lies on the
%management of measurements. As we have seen throughout the paper, there are
%multiple software components that need to be measured, which may come from
%different stakeholders. Besides, some of them undergo frequent updates that
%further increase the number of "golden measurements" to manage. Here, it may be
%desirable to rely on a trusted third party that keeps a database of acceptable
%measurements and associated metadata.

Our evaluation did not focus on userspace software that runs after booting the
root filesystem, such as \texttt{cloud-init}~\cite{cloudInit} or
\ac{CSP}-managed software like Azure Linux VM Agent~\cite{azureDocsLinuxAgent}.
These agents take care of VM configuration and can execute privileged commands,
like installing packages or updating the trusted SSH identities. Any such
operation, if compromised, can thwart the security guarantees of
\acp{CVM}. Hence, it is equally important to implement hardening measures in the
root filesystem to, e.g., check \texttt{cloud-init} configuration files before
executing them or disable unnecessary software.

Lately, several companies like Anjuna~\cite{anjuna}, Edgeless
Systems~\cite{edgelessSystems}, and Enclaive~\cite{enclaive} have started
offering their own confidential computing solutions to customers, providing
fully-managed confidential workloads in form of containers, virtual machines, or
both, and supporting deployments on popular \acp{CSP}. Attestation of workloads
is managed by a dedicated attestation service (which can also run within a
\ac{TEE}) and is therefore transparent to the customer, who only receives a
report indicating whether or not their deployment is verified. Such solutions
can largely improve user experience and ease the adoption of confidential
computing. However, when it comes to cloud deployments, they still depend on the
\acp{CSP} and the confidential computing features offered by them.  Therefore,
besides adding another third party to the trust model, such solutions cannot
increase the \emph{trustworthy AL} shown in \cref{tbl:cloud-landscape}.
Analogously, the same applies to third-party attestation services such as Intel
Trust Authority~\cite{intelTrustAuthority}.

A possible solution to the trust problem could leverage the live migration
feature of \acp{TEE} to initialize a \ac{CVM} locally and then migrate its state
to a remote host, where it continues its execution. Live migration is performed
over a secure channel established between the \ac{TEE} modules in both hosts,
preserving the confidentiality and integrity of the \ac{CVM}. Since the \ac{CVM}
is initialized on a local platform, the guest owner has full control over its
boot process and could therefore implement a full \emph{AL4} attestation
strategy. Live \ac{CVM} migration is supported by AMD \sevsnp{} and Intel
\ac{TDX}~\cite{AmdSevSnpWhitepaper,intelTdxMigration}. However, apart from some
technical and organizational challenges (e.g., the guest owner needs proper
\ac{TEE} hardware to initialize the \ac{CVM}), this workflow is currently not
supported by any \ac{CSP} and thus we did not explore it further.

Finally, we note that CC is still a relatively new technology and support for
recent \acp{TEE} such as \sevsnp{} in the cloud has not yet reached maturity. On
the evaluated \acp{CSP}, \acp{CVM} are currently available in selected regions
only and, on \ac{GCP}, support is in public preview. Despite existing
limitations, we see many efforts from \acp{CSP} to improve their offering,
driven in part by the open source community and standardization bodies. For
example, Google has recently open-sourced an unofficial tool to verify the
firmware of \acp{CVM} running on \ac{GCP}~\cite{gceTcbVerifier}, which
potentially allows reaching a trustworthy \emph{AL2}.
